\documentclass[12pt,a4paper]{article}
\usepackage[utf8]{inputenc}
\usepackage[vietnamese]{babel}
\usepackage[left=2cm,right=2cm,top=2cm,bottom=2cm]{geometry}
\usepackage{mathptmx}
\usepackage{enumitem}

\pagestyle{plain}
\linespread{1.15}

\begin{document}

% --- HEADER QUỐC HIỆU ---
\noindent
\begin{minipage}[t]{0.45\textwidth}
    \begin{center}
        \small
        BỘ CÔNG AN\\
        \textbf{\underline{VIỆN KHOA HỌC HÌNH SỰ}}\\[0.1cm]
        Số: 09/KL-KTHS\\
        \textit{V/v kết luận giám định tài liệu}
    \end{center}
\end{minipage}
\hfill
\begin{minipage}[t]{0.52\textwidth}
    \begin{center}
        \small
        \textbf{CỘNG HÒA XÃ HỘI CHỦ NGHĨA VIỆT NAM}\\
        \textbf{\underline{Độc lập - Tự do - Hạnh phúc}}\\[0.1cm]
        \textit{Hà Nội, ngày 25 tháng 07 năm 2026}
    \end{center}
\end{minipage}

\vspace{0.6cm}

\begin{center}
    \textbf{\large KẾT LUẬN GIÁM ĐỊNH KỸ THUẬT HÌNH SỰ}\\
    \textit{(V/v Giám định tính thật giả của tờ di chúc ký hiệu f2-1)}
\end{center}

\vspace{0.3cm}

\noindent Viện Khoa học Hình sự Bộ Công an kết luận về mẫu giám định di chúc gửi kèm Quyết định trưng cầu số 14/QĐ-CSHS như sau:

\section*{I. NỘI DUNG YÊU CẦU GIÁM ĐỊNH}
Giám định tính đồng nhất của nét chữ, chất liệu mực và thời điểm viết đoạn văn bản: *"thuộc Mã thửa địa chính khóa 2021-BS14"* trong tờ di chúc lập ngày 15/04/2018.

\section*{II. KẾT QUẢ GIÁM ĐỊNH}
\begin{enumerate}[label=\arabic*., leftmargin=1.5cm]
    \item \textbf{Phân tích nét chữ:} Đoạn văn bản chỉ định có đặc điểm nét chữ bị run, ngắt quãng bất thường, nét mực đè lên dòng chữ cũ của cụ Nguyễn Văn Thành.
    \item \textbf{Phân tích phổ quang quang học:} Đoạn văn bản chèn đè được viết bằng loại mực bi hóa dầu nhãn hiệu hiện đại sản xuất năm 2024, hoàn toàn khác biệt với chất liệu mực bi mực nước lập năm 2018.
\end{enumerate}

\section*{III. KẾT LUẬN GIÁM ĐỊNH}
Đoạn văn bản *"thuộc Mã thửa địa chính khóa 2021-BS14"* trong tờ di chúc bị **VIẾT CHÈN ĐÈ GIẢ MẠO** vào khoảng đầu năm 2024. Chữ ký ban đầu của cụ Thành bị cố tình sửa chữa nét để thâu tóm tài sản thừa kế.

\vspace{0.8cm}

\noindent
\begin{minipage}[t]{0.45\textwidth}
    \small
    \textbf{\textit{Nơi nhận:}}\\
    - Cơ quan CSĐT Công an TP;\\
    - Lưu: VT, CSHS.
\end{minipage}
\hfill
\begin{minipage}[t]{0.5\textwidth}
    \begin{center}
        \small
        \textbf{VIỆN TRƯỞNG VIỆN KHHS}\\
        \textit{(Ký và đóng dấu)}\\[1.5cm]
        \textbf{Thiếu tướng Trần Quốc Tuấn}
    \end{center}
\end{minipage}

\end{document}
